\section{The embedding problem}
\label{sec:introduction}

Throughout, $C^*$-algebras are complex, and an \emph{embedding} is an
injective $*$-homomorphism. We state unit preservation explicitly.
For a nonzero unital $C^*$-algebra $B$ and a free ultrafilter $\omega$ on
$\N$, its \emph{norm ultrapower} is
\[
 B^\omega=\ell^\infty(\N,B)/I_\omega(B),\qquad
 I_\omega(B)=\{(b_n):\lim_\omega\|b_n\|=0\}.
\]
Here $\ell^\infty(\N,B)$ is the algebra of bounded sequences, and the
quotient uses the operator norm. The Cuntz algebra $\Otwo$ is the universal
unital $C^*$-algebra generated by two isometries $s_1,s_2$ satisfying
$s_1s_1^*+s_2s_2^*=1$; it was introduced by Cuntz~\cite{Cuntz1977}.

The embedding theorem of Kirchberg--Phillips~\cite{KirchbergPhillips2000}
states that every separable unital exact $C^*$-algebra embeds unitally
into $\Otwo$. Here exactness means that minimal tensoring preserves
short exact sequences. Passing from $\Otwo$ to its norm ultrapower
asks how far this universality extends beyond exact algebras. The
techniques behind the embedding theorem led Kirchberg to consider this
possibility, as recounted in \cite[Section~1]{GoldbringSinclair2015}.

Kirchberg's embedding problem asks whether every separable unital
$C^*$-algebra admits a unital embedding into $\Otwo^\omega$; see
Goldbring--Sinclair~\cite[Section~3.1]{GoldbringSinclair2015}. We give an explicit
counterexample. The unital formulation is equivalent to the formulation
without unit preservation~\cite[Section~3.1]{GoldbringSinclair2015}.
This question concerns norm ultrapowers of a nuclear
$C^*$-algebra. It is distinct from the tensor-norm formulation of
Kirchberg's conjecture and from Connes' embedding problem, as emphasized
in~\cite[Section~1]{GoldbringSinclair2015}.

\Needspace{14\baselineskip}
\begin{theorem}\label{thm:main}
Let
\[
 G=\Z[1/2]^3\rtimes\bigl(\SL_3(\Z)\times\Z\bigr),
 \qquad (M,k)\cdot v=2^kMv,
\]
and let $A=C^*(G)$ be its full group $C^*$-algebra. Then $A$ is separable
and unital. For every nonzero unital nuclear $C^*$-algebra $B$ and every
free ultrafilter $\omega$ on $\N$, there is no unital embedding
\[
 A\longrightarrow B^\omega.
\]
In particular, $A$ does not embed unitally into $\Otwo^\omega$.
\end{theorem}

The group is countable and independent of both $B$ and $\omega$;
no separability assumption on $B$ is needed.
The proof treats an arbitrary fixed nuclear $B$ throughout. Specializing
to the nuclear algebra $\Otwo$ gives the negative answer to the embedding
problem; for its nuclearity, see~\cite[Section~4, pp.~80--82]{PaulsenZheng2016}.
Section~\ref{sec:consequences} gives the consequences for existential
closure, permanence of $\Otwo^\omega$-embeddability under crossed products,
and good nuclear witnesses, with the relevant notions defined there.

\subsection{Previous work and the new obstruction}

\begingroup
\setlength{\emergencystretch}{2em}
Recent approaches study computability and tensor norms.
Fox--Goldbring--Hart~\cite{FoxGoldbringHart2024} derive
computability-theoretic consequences of a positive answer, using the
computable presentation of $\Otwo$. Goldbring--Sinclair give a conditional
route to a negative answer through a \emph{weakly explicit minimal tensor
norm}: an effective choice of universal formulas approximating the
minimal tensor-norm predicates to a fixed accuracy.
They leave this property of $\Otwo$ open
\cite[Definition~4.16, Theorem~4.18 and Question~4.20]{GoldbringSinclair2025Tensor}.
Our obstruction instead comes from the representation theory of an
explicit discrete group.\par
\endgroup

The group configuration already appears in the nonfiniteness arguments of
Sauers~\cite[Section~2.4]{Sauers2026NonMF} and
Eckhardt~\cite[Theorem~3.1 and Example~3.2]{Eckhardt2026}.
They show that the full algebra of the ascending HNN extension $G$ is
nonfinite: it contains a proper isometry $v$, so $v^*v=1$ but $vv^*\ne1$.
Its property~(T) base group
$\Gamma=\Z^3\rtimes\SL_3(\Z)$ and its finite-index scalar
self-embeddings are recorded in
de Cornulier~\cite[Proposition~1.1 and its proof]{deCornulier2007}.
Sauers and Eckhardt also construct non-MF groups: groups admitting no injective
homomorphism into the unitary group of any norm matrix corona
$\prod_j M_{d_j}/\bigoplus_j M_{d_j}$, with positive integer dimensions
$d_j$. Here the product consists of bounded sequences and the direct sum
is the ideal of sequences tending to zero in norm
\cite[Theorem~4.3]{Eckhardt2026}\cite[Sections~3--5]{Sauers2026NonMF}.
Eckhardt's example is an associated wreath product. The target class here is different:
nonfiniteness of $C^*(G)$ alone cannot exclude a unital embedding into
$\Otwo^\omega$, which itself contains proper isometries.

The additional ingredient is a finiteness statement for approximate
finite-dimensional representation types inside a nuclear algebra
(Lemma~\ref{lem:nuclear-finiteness}). Its tensor-product argument follows
Wassermann's method~\cite{Wassermann1991}, as explained
by Manuilov--Thomsen~\cite[Section~2.1]{ManuilovThomsen2007}.
The statement applies to an arbitrary tuple of unitaries, even when the
tuple satisfies no group relations. This flexibility is essential because
coordinate lifts of an ultrapower representation are only approximately
multiplicative. A dimension-independent estimate allows the representation
type to vary with the coordinate. Finite-index induction then relates
members of the resulting finite sets by composition with the self-embedding.
These two passages turn the group configuration into the norm-ultrapower
obstruction.

\subsection{Proof outline}

The subgroup $\Gamma=\Z^3\rtimes\SL_3(\Z)$ of $G$ has property~(T).
Conjugation by the generator $t$ of the last $\Z$-factor sends $\Gamma$
onto the index-eight subgroup
$\theta(\Gamma)$, where $\theta(v,M)=(2v,M)$.
There is a finite-dimensional irreducible representation $\pi_*$ of
$\Gamma$ such that $\pi_*\circ\theta^j$ is inequivalent to $\pi_*$ for
every $j\ge1$. Section~\ref{sec:group} verifies these facts.

Fix a finite generating set $S$ of $\Gamma$ for which property~(T)
gives a spectral gap. Suppose a unital embedding $A\to B^\omega$ exists,
and choose unitary lifts $x_n(g)\in B$ of its group unitaries. For each
$n$, consider the nonzero finite-dimensional irreducible unitary representations
$\pi$ for which the tuple $(x_n(s)\otimes\bar\pi(s))_{s\in S}$ has
spectral minimum below a fixed small threshold for the sum of squared
displacements from~$1$. Here $\bar\pi$ is the conjugate representation.
Sections~\ref{sec:spectral} and~\ref{sec:nuclear} specify the threshold
and show that the resulting set $D_n$ of representation classes is finite.

Write $\pi\prec\rho$ when $\pi$ occurs as a unitary subrepresentation
of $\rho$, and call $\pi'$ a predecessor of $\pi$ if
$\pi\prec\pi'\circ\theta$. Faithfulness of the embedding places
$[\pi_*]$ in $D_n$ on a set belonging to $\omega$.
In Section~\ref{sec:contradiction}, the spectral gap for
\emph{every sequence} of finite-dimensional representations supplies,
on one set belonging to $\omega$, the simultaneous implication
\[
 [\pi]\in D_n\quad\Longrightarrow\quad
 \text{some }[\pi']\in D_n\text{ satisfies }\pi\prec\pi'\circ\theta.
\]
Starting at $\pi_*$ and repeatedly choosing a predecessor must give a
repetition in the finite set $D_n$. Irreducibility then makes $\pi_*$
periodic under composition with $\theta$, a contradiction.
